%!TEX root = ../thesis.tex
%*******************************************************************************
%****************************** Eight Chapter **********************************
%*******************************************************************************
% **************************** Define Graphics Path **************************
\ifpdf
    \graphicspath{{content/chapter9/figs/raster/}{content/chapter9/figs/pdf/}{content/chapter9/figs/}}
\else
    \graphicspath{{content/chapter9/figs/vector/}{content/chapter9/figs/}}
\fi

\chapter{Conclusion}\label{chapter:conclusion}

% PART 1: Recap of the Thesis Aims and Objectives. Purpose: Begin by restating the primary goals outlined in your Motivation and Introduction chapters.

The overarching aim of this thesis has been to address a central challenge in modern environmental science: extracting meaningful insights from an ever-growing array of environmental data sources. Over recent decades, the volume and diversity of these sources---spanning dense, regularly sampled climate reanalysis to sparse, irregular observations of marine litter---has expanded at a pace requiring advanced statistical and computational solutions. In particular, the proliferation of high-resolution observations calls for efficient dimensionality reduction techniques, while data-poor scenarios demand robust Bayesian frameworks for dealing with uncertainty and uneven sampling.

Building on these motivations, the thesis combines explorative and hypothesis-driven approaches to illuminate patterns in both dense and sparse data regimes. For dense climate datasets, the focus was on refining linear decomposition techniques to capture dominant variability modes and lagged teleconnections, including seasonal effects. Meanwhile, for increasingly concerning issues like marine litter, where observations are often sporadic and inconsistent, a Bayesian Gaussian process regression framework was introduced to test hypotheses regarding the seasonality and principal drivers of macroplastic pollution. 

% PART II: Summary of Main Findings. Purpose: Provide a concise synthesis of what was accomplished, integrating results from Sections xeofs, tele, and beach_litter without duplicating their individual conclusions.

A core goal of this thesis was to demonstrate how advanced dimensionality reduction and Bayesian machine learning methods can tackle distinct challenges posed by dense and sparse environmental datasets. Three main strands of research were explored:

\begin{enumerate}
    \item \textbf{Efficient and Scalable Dimensionality Reduction (Chapter~\ref{chapter:xeofs})}\\To address the increasing volume and complexity of climate data, this work introduced the \software{xeofs} library. By leveraging the Python-based Pangeo ecosystem (notably \software{xarray} and \software{dask}), \software{xeofs} facilitates both traditional and advanced dimensionality reduction on large, multidimensional datasets. Through randomized linear algebra, it achieves order-of-magnitude improvements in speed compared to existing tools. Its open-source, modular design also provides a flexible platform for future expansions, such as incorporating nonlinear algorithms or more sophisticated bootstrapping methods. Together, these enhancements enable researchers to more readily apply dimensionality reduction techniques, from personal workstations to HPC clusters, thereby bridging the gap between theoretical methods and practical climate data analysis.
    \item \textbf{Extracting Shared Variability and Lagged Teleconnections (Chapter~\ref{chapter:tele})}\\Turning to dense reanalysis data, the thesis applied a rotated Hilbert Maximum Covariance Analysis (MCA) framework to uncover key modes of shared ocean-atmosphere variability in global SST and continental precipitation. This approach successfully isolated known climate phenomena---such as the seasonal cycle, ENSO, and the PDO---demonstrating the method's robustness. Additionally, it captured new or less-studied teleconnections, including subtropical-to-tropical links and possible Arctic influences. By operating in a complexified, phase-sensitive domain, the analysis could simultaneously resolve seasonal oscillations and lagged interactions without removing the seasonal cycle a priori. Complementary correlation-based extensions (CPCCA) yielded alternative perspectives on climate variability patterns. This highlights the dual utility of these methods both as exploratory tools and as a basis for generating new, testable hypotheses about large-scale climate interactions.
    \item \textbf{Modeling Seasonal Patterns in Sparse Data (Chapter~\ref{chapter:beach_litter})}\\In scenarios where data are irregular and scattered, such as in beach litter observations, this thesis demonstrated the utility of log-Gaussian Cox process (LGCP) regression within a Bayesian framework. The analysis revealed robust seasonal macroplastic patterns along certain coastal regions of the North-East Atlantic and identified key drivers—particularly small-scale rivers and aquaculture activities—in shaping these patterns. Crucially, the Bayesian approach allowed for thorough uncertainty quantification, underscoring the importance of sufficient temporal resolution in monitoring strategies. These findings also stressed the need to refine existing litter classification protocols and to focus on root causes, such as plastic emissions from mariculture. By combining detailed statistical models with targeted domain knowledge, the work illuminates how even limited observational data can yield important environmental insights when approached with robust statistical and computational methods.
\end{enumerate}

%%%%% Reflections on Methods and Contributions: Purpose: Critically assess the strengths, limitations, and significance of the approaches used: %%%%


This thesis employed two main methodological pillars---rotated Hilbert MCA/CPCCA and Bayesian LGCP regression---to tackle the diverse challenges posed by dense and sparse environmental datasets. Each approach brought distinct advantages but also encountered notable limitations, opening up possible avenues for further refinements.

\paragraph{Rotated Hilbert MCA/CPCCA}
These linear dimensionality-reduction methods proved invaluable for extracting dominant patterns of shared variability in dense climate datasets. By integrating Hilbert transforms, the framework provided amplitude-phase representations that facilitate the identification of regularly oscillating modes, such as seasonal cycles or time-lagged teleconnections. Additionally, randomization-based linear algebra scaled effectively with the large number of observations and variables often present in reanalysis data. However, MCA/CPCCA relies on strong linear assumptions, potentially missing the richer complexity of nonlinear phenomena like ENSO's multiple feedback processes. The interpretation of amplitude-phase information also loses clarity for irregular or multi-frequency signals. Moreover, covariance-based methods alone do not yield reduced-order dynamical models, and their purely frequentist underpinnings mean that comprehensive uncertainty assessments require computationally expensive bootstrapping. Despite these constraints, MCA/CPCCA excel as exploratory tools for dense datasets, illuminating previously unrecognized teleconnections and guiding subsequent, more targeted investigations.

\paragraph{Bayesian LGCP regression}
In contexts where data are irregular or scarce, LGCP's Bayesian foundations offer crucial advantages. By explicitly modeling the covariance structure through kernels, LGCP regression handles complex spatial and temporal dependencies while providing direct estimates of posterior uncertainty. This made it particularly appropriate for modeling macroplastic beach litter data, where monitoring efforts yield sporadic coverage and missing observations. That said, LGCP regression can rapidly become computationally expensive, necessitating careful approximations or specialized libraries for large-scale deployments. Even so, its capacity to merge different data streams (e.g., citizen science and operational measurements) suggests wide-ranging applicability to environmental problems where data scarcity impedes simpler methods. Moreover, GPR's built-in uncertainty quantification helps guide more effective resource allocation for monitoring, as demonstrated by identifying seasonal hotspots and potential anthropogenic drivers in beach pollution data.

Taken together, these two methodological pillars illustrate the trade-offs inherent in environmental data analysis. MCA/CPCCA excel in high-dimensional, data-rich scenarios by revealing large-scale patterns and teleconnections, albeit with linear constraints. GPR's Bayesian structure, meanwhile, shines in uncertain, data-poor conditions but faces scaling issues if the dataset grows too large. Their collective contribution demonstrates that no single technique suffices for all environmental contexts; rather, careful selection and tailoring of methods, aided by supportive open-source software, can address the increasingly complex challenges at the intersection of big data and Earth system science.

% Comparison to Existing Literature. Purpose: Place your findings in the broader context of ongoing research.

The approaches presented in this thesis align broadly with ongoing developments in climate data science and environmental modeling. The rotated Hilbert CPCCA framework, for instance, builds on earlier work by \citet{swenson_continuum_2015} by combining the CPCCA framework with Hilbert transform and rotation, concepts borrowed from the PCA literature, allowing us to can untangle phase-resolved oscillatory modes, similar to parallel efforts in the PCA literature by \citet[e.g.][]{bueso_nonlinear_2020,guilloteau_rotated_2020}.  Likewise, the detection of lesser-known or multifrequency teleconnections complements prior findings on canonical ENSO, PDO, and NAO modes \cite[e.g.][]{timmermann_ninosouthern_2018, mezzina_dynamics_2020,wiedermann_differential_2021,tang_positive_2024}, reaffirming the utility of MCA-like methods as exploratory tools in climate diagnostics.

The development of \software{xeofs} echoes the trend toward specialized software libraries that integrate seamlessly with distributed computing frameworks, such as \software{xarray} and \software{dask}, and the broader Pangeo ecosystem \cite{pangeo}. By leveraging randomized linear algebra, the present study refines these techniques for large-scale datasets, paralleling other recent efforts to optimize PCA-based analyses in multi-terabyte datasets \cite[e.g.][]{velegar_scalable_2019,erichson_sparse_2020}. Its emphasis on open-source collaboration and modular design reflects a broader community-driven approach, seen in projects like \software{xarray} and \software{dask}, and it resonates with the FAIR data principles \cite{wilkinson_fair_2016} adopted in various Earth science consortia.

Meanwhile, applying a Bayesian LGCP to sparse macroplastic litter data extends a growing body of research on data assimilation, citizen science, and risk assessment in environmental pollution studies \cite[e.g.][]{datta_nonseparable_2016,thiel_impacts_2018,locritani_assessing_2019,chenet_plastic_2021,zorzo_approach_2021,haarr_global_2022}. Previous analyses often relied on limited local or spatially aggregated data, whereas this spatially resolved modeling---enhanced by explicit uncertainty quantification---demonstrates the value of Bayesian methods in handling point processes and merging diverse data streams. Although earlier studies emphasized riverine inputs and fishing as the major anthropogenic sources of beach litter, our findings reveal that small-scale rivers and, more notably, mariculture-related items may serve as key drivers of seasonal variability in pollution hotspots. This observation aligns with the growing recognition of litter seasonality and aquaculture-driven plastic contamination \cite{skirtun_plastic_2022,ruangpanupan_seasonal_2023,xu_seasonality_2024,cozar_proof_2024} and strengthens calls for strategic monitoring, targeted cleanup efforts, and adaptive policy-making in marine litter mitigation, as proposed in recent policy-focused analyses \cite{ospar_comission_feeder_2021, hipolito_policy_2020}.

Overall, these results situate the thesis contributions within the broader evolution of Earth system analytics, which is propelled by increasingly large, multi-source datasets. The combination of advanced dimensionality reduction and Bayesian machine learning echoes contemporary research directions emphasizing both powerful explorative methodologies for dense datasets and rigorous, uncertainty-aware frameworks for sparse or noisy settings.

% "Limitations and Critical Evaluation". Purpose: Acknowledge the scope and constraints of your thesis to maintain academic rigor.

% Future Directions

%     Purpose: Propose how subsequent studies might extend or build upon your work.
%     Key Points:
%         Potential enhancements to the applied methods (e.g., more sophisticated Bayesian frameworks, deep learning approaches, advanced HPC solutions).
%         Exploration of additional datasets or new problem domains (e.g., microplastics, other marine litter categories, or teleconnections in other parts of the world).
%         Strategies to integrate advanced visualization or interpretability techniques, bridging domain science and machine learning further.

\section*{Future Directions}
Several promising lines of inquiry and methodological enhancements emerge from the findings and methods presented in this thesis. One important avenue lies in the further exploration of climate teleconnection patterns, particularly in the hydrological cycle. By incorporating additional physical variables such as moisture flux or sea surface salinity, it may be possible to refine the understanding of how large-scale correlations influence sub-seasonal to decadal variability \cite[e.g.][]{ballabrerapoy_potential_2002,rathore_improving_2021,li_skillful_2022}. A deeper knowledge of these correlations could also lead to improved climate forecasts, especially if teleconnections are integrated into operational models that focus on precipitation patterns or marine conditions \cite{li_skillful_2022,straaten_correcting_2023}.

In addition to refining the teleconnection analyses, there is considerable scope for using the \software{xeofs} software library to investigate large, multidimensional environmental datasets outside the scope of the examples provided in this thesis. By systematically applying its available dimensionality reduction methods to new datasets, researchers can both compare different algorithms and tailor the choice of method to specific scientific questions. Indeed, systematically performing such comparisons could shed light on how linear decomposition techniques fare in varying contexts, potentially helping to clarify which approaches are best suited for specific data characteristics or research goals.

A further direction involves deepening the Bayesian modeling strategy that proved successful in analyzing spatio-temporal macroplastic distributions. Building on this success, future studies might integrate multiple data sources, such as citizen science observations and high-accuracy government or satellite measurements, into a coherent Bayesian framework. Such an approach could resonate with broader initiatives, such as the Global Litter Observatory project \cite{marine_litter_lab_global_2021}, that seek to unify scattered litter-monitoring datasets and offer robust spatio-temporal insights into plastic pollution. In addition, adapting Gaussian process regression for classification tasks, for instance in predicting marine species occurrences using the recently developed MINKA database \cite{luna_reyes_evaluation_2022}, opens opportunities to serve a variety of ecological and conservation-related questions. 

For dimensionality reduction specifically tailored to sparse regimes, advancing beyond the frequentist methods used here, a promising extension lies in deploying Bayesian regularization to prevent overfitting when data coverage is partial or highly uneven. Approaches like Gaussian-process factor analysis \cite{yu_gaussian-process_2008,luttinen_variational_2009}, which combines factor analysis with a Gaussian-process kernel, could provide exploratory capabilities akin to PCA, yet offer uncertainty quantification essential for drawing reliable inferences from incomplete datasets. The application of such methods to pan-European beach litter observations \cite{molina_jack_emodnet_2019}, for example, might uncover subtle trends or spatial distributions that elude standard frequentist techniques.



% Concluding Remarks

%     Purpose: Offer a concise, final perspective on how the thesis contributes to the broader scientific conversation.
%     Key Points:
%         Re-emphasize the central theme: bridging dense and sparse data regimes via machine learning and dimensionality reduction.
%         End with an optimistic note on how these findings enhance understanding of the Earth system and encourage more data-driven decision-making.

Overall, this thesis illustrates how advanced dimensionality reduction and Bayesian modeling can enhance our capacity to extract meaningful insights from environmental datasets, whether dense or sparse. By demonstrating the utility of linear explorative methods such as rotated Hilbert MCA/CPCCA for uncovering shared variability, introducing \software{xeofs} to streamline and scale dimensionality reduction, and applying Gaussian process regression in data-scarce contexts, it bridges theoretical innovations with practical needs in climate and environmental research. The ability to handle increasingly large and varied data sources highlights both the urgency and the potential for continued methodological refinement. Ultimately, these contributions underscore the importance of robust, flexible, and open-source approaches in our collective endeavor to better understand Earth system processes, address pressing environmental challenges, and inform future policy and resource management decisions.



%%%%% Reflections on Methods and Contributions: Purpose: Critically assess the strengths, limitations, and significance of the approaches used: %%%%

% rotated Hilbert MCA strengths:
% - can identify dominant patterns of variability in large climate data sets, such as teleconnections
% - linear assumption makes mathematical interpretation of the patterns mostly straightforward
% - amplitude-phase information can be used to identify regularly oscillating patterns in the data
% - when implemented using ranomdized linear algreba it scales well with the number of observations and variables, making it suitable for high-dimensional data sets, such as climate data.

% MCA limitations:
% - amplitude-phase information is less interpretable for irregular oscillatory modes
% - non-linear phenomena are not well captured by the linear assumption
% - interpretation of the patterns is based on the covariance/correlation structure in the data, which may not be sufficient to derive reduced-order dynamical models
% - non-bayesian approach, uncertainty quantification requires bootstrapping which may become unfeasible for large data sets; not directly applicable to sparse data sets with many missing values

% significance of the approach:
% - great explorative tool for the analysis of dense data sets, that allows to identify patterns shard patterns between different variables that are not immediately obvious from prior knowledge


% xeofs strengths:
% - easy to use and well-documented software package that allows to perform a variety of popular dimensionality reduction techniques on dense data sets
% - seamless integration with xarray and dask makes it well-suited for the analysis of large data sets, including but not limited to climate data specifically
% - achieves an acceleration of about one order of magnitude in computational time compared to comparable software packages
% - embedded into the state-of-the-art Python-based Pangeo software stack, making it accessible to a broad audience
% - open-source philosophy allows for easy reporting, analyzing, and solving of errors and bugs

% xeofs weaknesses:
% - limited scope to linear dimensionality reduction methods
% - lack of support for some traditional methods and recent innovations
% - offerns only non-bayesian frameworks


% xeofs significance:
% - step in the right direction to make dimensionality reduction techniques more accessible to the broader scientific community
% - contributes to the ever-growing challenge to make sense of the vast amount of environmental data generated by satellites, sensors, and field observations
% - first signs emerge that xeofs has the potential to become a widely used software package in the field of climate science and beyond. it has recently been included in the closer xarray ecosystem and is part of the Pangeo software stack featured through its education working group, the Pythia project.



% GPR strengths:
% - Bayesian framework allows for uncertainty quantification and regularization of the model
% - therefore can model complex spatio-temporal patterns in the data, even when data is scarce and irregularly sampled
% - can be used to model the joint spatial distribution of different data sources, making it ideal for combining citizen data with high-precision government data, offering an avenue for many follow-up studies

% GPR weaknesses:
% - computationally expensive, especially for large data sets, however can be mitigated by using sparse approximations


% GPR significance:
% - Bayesian approaches offer a powerful approach in noisy scarce data regimes, where overfitting becomes a major concern. spatio-temporal distribution modeling of environmental pollution such as marine litter, is a prime example and was demonstrated to work well here.
% - modeling the covariance structure of different samples implicitly by a kernel function, GPR becomes extremley flexible and can be used in a variety of environmental problems, where we have sparse data sets and even a different variety of data sources.
% - lots of potential in assimilation of different data sources, such as citizen data and high-precision government data, to model the joint spatial distribution based on the different underlying accuracy of different data sources
% - great potential as a environmental monitoring tool, that may guide monitoring strategies, such as seasonal cleanup efforts, and inform marine litter mitigation strategies in the future. interest by the spanish ministry has been expressed and a current collaborative follow-up study is currently underway with representatives from the spanish ministry of environment. 
